%% Complete proofs for the active primitive-language Tseitin construction.
\begin{annexitem}{annex:p2.prop.tseitin}{thm:p02-tseitin}{Definitional extension and equisatisfiability}
\begin{proof}
Use the notation of Definition~\ref{def:p02-tseitin-construction}.
Unique readability and the finite occurrence tree make every representative
and gate block well defined.  A finite valuation extends to all of \(\mathsf{Prop}\)
by assigning \(\mathbf{F}\) elsewhere; the coincidence lemma makes evaluation
on its supported formulas independent of those other values.
Lemma~\ref{lem:prop-clause-satisfaction} says that a clause is satisfied
exactly when at least one of its literals is true, and that the conjunction
of clause interpretations is satisfied exactly when all clauses are.

\paragraph{Every primitive gate.}
At a variable occurrence \(\varphi_a=p_n\), the representative is \(p_n\)
itself and there is no constraint to check.  At a \(\bottom\) occurrence,
the unit clause \(\{\neg z\}\) holds exactly when \(w(z)=\mathbf{F}\), the
required value of the nullary connective.

At a negation node the clauses are \(\{\neg z,\neg x\}\) and
\(\{z,x\}\).  If \(w(z)=\mathbf{T}\), the first forces \(w(x)=\mathbf{F}\)
and the second is satisfied.  If \(w(z)=\mathbf{F}\), the second forces
\(w(x)=\mathbf{T}\) and the first is satisfied.  Conversely each of these two
allowed pairs satisfies both clauses.  Thus the clauses hold exactly when
\(w(z)\) is the negation of \(w(x)\).

At a conjunction node, if \(w(z)=\mathbf{T}\), the first two clauses force
both \(w(x)=\mathbf{T}\) and \(w(y)=\mathbf{T}\), and the third is satisfied.
If \(w(z)=\mathbf{F}\), the first two are satisfied and the third requires
at least one of \(w(x),w(y)\) to be false.  Conversely these conditions
satisfy all three clauses.  They say exactly that \(w(z)\) is the
conjunction of the two input values.

At a disjunction node, if \(w(z)=\mathbf{F}\), the last two clauses force
both \(w(x)=\mathbf{F}\) and \(w(y)=\mathbf{F}\), while the first is satisfied.
If \(w(z)=\mathbf{T}\), the last two are satisfied and the first requires
at least one input to be true.  Conversely these conditions satisfy all
three clauses.  They say exactly that \(w(z)\) is the disjunction of the
two input values.

At an implication node, if \(w(z)=\mathbf{F}\), the clauses \(\{z,x\}\)
and \(\{z,\neg y\}\) force \(w(x)=\mathbf{T}\), \(w(y)=\mathbf{F}\),
and the first clause is satisfied.  If \(w(z)=\mathbf{T}\), those two clauses
are satisfied and the first requires \(w(x)=\mathbf{F}\) or
\(w(y)=\mathbf{T}\).  Conversely these conditions satisfy all three clauses.
They are exactly the falsity and truth conditions for implication.
These arguments do not require \(x\) and \(y\) to be different variables.
They check every primitive constructor and both directions of each gate
definition.

\paragraph{Existence of the defining extension.}
Let \(v:V\to\mathbb{B}\).  Define \(w\) on \(V\cup Z_\varphi\) by
\[
 w(p)=v(p)\quad(p\in V),\qquad
 w(z_a)=\widehat v(\varphi_a)\quad(a\in G_\varphi).
\]
Freshness and distinctness make this a well-defined function extending
\(v\), even when \(V\) or \(Z_\varphi\) is empty.
This definition evaluates original subformulas under \(v\); it is not a
circular system of assignments to the new variables.
At a variable occurrence,
\(w(r_a)=v(p_n)=\widehat v(\varphi_a)\); at every other occurrence the
same equality is the definition of \(w(z_a)\).
Consequently the representative values at every node and its children
satisfy the appropriate truth function.  The gate verification above gives
\(w\entails\mathcal D_\varphi\).

\paragraph{Every defining extension has these values.}
Conversely, let \(w:V\cup Z_\varphi\to\mathbb{B}\) extend \(v\) and satisfy
\(\mathcal D_\varphi\).  Induct from the leaves upward in the original
occurrence tree.  At a variable leaf the required identity follows from
\(w|_V=v\).  At a \(\bottom\) leaf the unit clause forces the value
\(\mathbf{F}=\widehat v(\bottom)\).
At a negation node, the induction hypothesis identifies the child's
representative value with its value under \(v\).  The two negation clauses
therefore force the representative at the node to have the value of the
original negation.  At a conjunction node the two child identities and the
conjunction gate give the conjunction value.  At a disjunction node those
identities and the disjunction gate give the disjunction value.  At an
implication node those identities and the implication gate give the
implication value.  All nodes have now been covered, proving
\(w(r_a)=\widehat v(\varphi_a)\) for every position.

In particular, every new variable's value is forced, and values on \(V\)
were fixed from the outset.  Thus the defining extension is unique on
\(V\cup Z_\varphi\).  This also handles an atomic input: there are no
auxiliary variables or defining clauses, and the unique extension is \(v\).

\paragraph{The root and the CNF formula.}
The defining extension satisfies the added clause \(\{r_\epsilon\}\)
exactly when \(w(r_\epsilon)=\mathbf{T}\).  By the identity just proved, this
is equivalent to \(\widehat v(\varphi)=\mathbf{T}\).  If the input is false
under \(v\), no extension satisfying the rooted output exists; if it is
true, the defining extension is its unique such extension.

The emitted list and its underlying set require precisely the same clauses
to be true: repetition changes neither condition.  The inherited clause
and conjunction semantics therefore give
\(w\entails T_\varphi\) iff \(w\entails S_\varphi\).
Taking existence over the finite input valuations and using their total
extensions to \(\mathsf{Prop}\) proves equisatisfiability of all three objects.
Finally, any total valuation satisfying \(T_\varphi\) restricts to a
model of \(S_\varphi\) on \(V\cup Z_\varphi\); its restriction to
\(V\) satisfies \(\varphi\).  Coincidence gives
\(\{T_\varphi\}\entails\varphi\) under total valuations as claimed.
The countervaluation in Example~\ref{def:p02-tseitin-root-boundary} shows why
the reverse entailment is not a general conclusion.
\end{proof}
\end{annexitem}

\begin{annexitem}{annex:p2.prop.tseitin-size}{thm:p02-tseitin-size}{Structural size and encoding}
\begin{proof}
Let \(n_0\) count \(\bottom\) occurrences, \(n_1\) negations, and
\(n_2\) binary-connective occurrences.  Then
\(M=n_0+n_1+n_2\); all other positions are original variable leaves.
The defining blocks and the root emit
\[
 n_0+2n_1+3n_2+1\leq3M+1
\]
clauses.  Before repeated literals within a clause are identified, their
total literal count is
\[
 n_0+4n_1+7n_2+1\leq7M+1.
\]
Set formation can only reduce these counts, and no listed clause has more
than three literals.  In \(\operatorname{form}_{\mathsf{CPL}}(C)\), each literal
uses at most one negation and each binary disjunction contributes one
connective and its parentheses.  Right-associated conjunction adds only a
constant number of tokens per clause.  Hence both the emitted list and its
CNF formula have \(O(N)\) token size under unit-cost names.

Preorder visits every tree node once.  With the children represented by stored
variable references and the fresh names already available, each node emits
a fixed bounded amount of data from its row.  Sorting at most three literals
per clause uses a bounded number of comparisons when names and comparisons
are unit-cost.  This proves the stated conditional \(O(N)\) emission bound;
it makes no assertion that arbitrarily long variable names cost one bit.

For the fresh-name bound, put \(d=|V|\).  There are \(N-M\) original
variable occurrences, so \(d\leq N-M\).  Of the first \(N\) symbols
\(p_0,\ldots,p_{N-1}\), at most \(d\) belong to \(V\).  At least
\(N-d\geq M\) are therefore available.  The first \(M\) unused symbols
chosen by the policy all have indices below \(N\).  This proves freshness
constructively and covers \(M=0\), as well as variable-free formulas with
\(d=0\).

Now let \(L\) be the stated bit length.  Each original variable leaf appears
in at most two emitted literals belonging to its parent block; for an atomic
input it appears once, in the root clause.  Thus copies of original variable
names contribute \(O(L)\) bits in total.  There are \(O(N)\) auxiliary
literal occurrences, each with an index requiring
\(O(\log(N+1))\) bits.  Connectives, polarity signs and separators add
\(O(N)\) bits.  The resulting output bound is
\(O(L+N\log(N+1))\).

An effective implementation first parses the finite input and computes
\(N\).  It then checks the finitely many candidate indices below \(N\)
against the input variables, assigns the required fresh indices in order,
visits the nodes, and writes the fixed clause blocks.  Even a direct scan of
the input for each candidate uses polynomially many bit comparisons; the
bounded-width clause sorting and output writing are polynomial as well.
No scan up to the largest original variable index is needed.  Position words
identify occurrences in the mathematical description; they are not used as
the binary output names.  Finally, \(N\) counts the actual primitive tree,
so repeated subformula occurrences and any expansion of defined connectives
are included rather than hidden in the size measure.
\end{proof}
\end{annexitem}
